\section{Conclusioni}
\subsection{Risultati dell'analisi}
Su npm il codice malevolo è entrato senza modificare il software dei pacchetti: un'e-mail di phishing ha permesso di raccogliere le credenziali di un manutentore reale e, con quelle, l'attaccante ha pubblicato nuove versioni che molti utenti hanno ricevuto come dipendenze transitive, spesso senza saperlo \cite{aikido2025,debug2025}. Su AUR chi attaccava ha adottato ricette rimaste senza manutentore, con un account proprio e del tutto legittimo, e ha aggiunto agli script d'installazione comandi che avviano codice esterno durante la transazione di \texttt{pacman} \cite{arch-aur-incident-2026,aur-mail-gnome-randr}. In entrambi i casi l'utente si è fidato di qualcuno che aveva, o sembrava avere, il diritto di pubblicare: su npm quel diritto era stato rubato, su AUR era stato ottenuto senza infrangere le regole della piattaforma.

Dal punto di vista del Cyber Risk Management ne segue che la procedura di adozione delle ricette orfane è di per sé una superficie d'attacco, anche quando nessun account viene compromesso. Una 2FA con chiave di sicurezza fisica avrebbe probabilmente fermato l'attacco a npm, perché una chiave FIDO è legata al dominio reale e non risponde a un sito contraffatto, come invece può accadere con un codice temporaneo inserito a mano \cite{npm-2fa-docs-2026}; sulle ricette AUR adottate non avrebbe avuto alcun effetto. Lì l'unico controllo previsto è la lettura di \texttt{PKGBUILD} e script \texttt{.install} da parte dell'utente \cite{arch-aur-incident-2026}, e non sappiamo quanti la facciano davvero a ogni aggiornamento. Il periodo di attesa prima di accettare una nuova release, come il cooldown di Dependabot o l'età minima di pnpm, è tra queste misure l'unica che lo sviluppatore imposta una volta sola nella configurazione \cite{github-disrupting-2026,pnpm-supply-chain}; protegge però solo se qualcuno, nel frattempo, segnala il problema e la versione viene ritirata. Nel caso npm la segnalazione è arrivata in tempi brevi, quindi anche un'attesa di pochi giorni sarebbe bastata \cite{aikido2025,debug2025}.

\subsection{Limiti e rischio residuo}
I circa due miliardi di download settimanali descrivono la diffusione dei pacchetti prima dell'attacco \cite{debug2025}: da questo dato non si ricava quante installazioni abbiano ricevuto le versioni malevole nel periodo in cui sono rimaste pubbliche, né quante siano finite in applicazioni usate da utenti reali. Il danno economico documentato è rimasto contenuto \cite{cso2025}. Le fonti su AUR spiegano come l'infostealer veniva installato e che cosa poteva fare \cite{arch-aur-incident-2026,ioctl2026}, ma non quanti sistemi abbia colpito. Data la disomogeneità dei dati, un punteggio numerico di rischio avrebbe dato un'idea di precisione non veritiera. Per lo stesso motivo \texttt{gnome-randr-rust} è usato solo come esempio: è il pacchetto su cui si trovano più informazioni, ma non si può dare per scontato che gli altri pacchetti dell'ondata funzionassero allo stesso modo.